 Ahora consideraremos la ecuación diferencial lineal homogénea de segundo orden escrita en forma estándar:
\[ y'' + p(x)y' + q(x)y = 0. \]

Nuestro objetivo es determinar una representación en serie de la solución general a esta ecuación diferencial centrada en $x = x_0$. Veremos que la existencia y forma de la solución dependen del comportamiento de las funciones $p$ y $q$ en $x = x_0$.

\begin{definition}
    El punto $x = x_0$ se llama un \textbf{punto ordinario} de la ecuación diferencial
\begin{equation} \tag{4.2.1}
y'' + p(x)y' + q(x)y = 0,
\end{equation}
si $p$ y $q$ son \textit{ambas} analíticas en $x = x_0$. Cualquier punto que no sea un punto ordinario de (4.2.1) se llama un \textbf{punto singular} de la ecuación diferencial.

\end{definition}

\begin{example}
    La ecuación diferencial
\[ y'' + \frac{1}{x^2 - 4}y' + \frac{1}{x + 1}y = 0 \]
tiene
\[ p(x) = \frac{1}{x^2 - 4} \quad \text{y} \quad q(x) = \frac{1}{x + 1}. \]

Vemos por inspección que los únicos puntos en los que $p$ falla en ser analítica son $x = \pm 2$, mientras que $q$ es analítica en todos los puntos excepto en $x = -1$. En consecuencia, los únicos puntos singulares de la ecuación diferencial son $x = -1, \pm 2$. Todos los demás puntos son puntos ordinarios. 
\end{example}

\begin{remark}
    En esta sección restringimos nuestra atención a los puntos ordinarios. Dado que las funciones $p$ y $q$ son analíticas en un punto ordinario, podríamos sospechar que cualquier solución de 4.2.1) válida en $x = x_0$ también es analítica allí y, por lo tanto, se puede representar como una serie de potencias convergente
\[ y(x) = \sum_{n=0}^{\infty} a_n(x - x_0)^n, \]
para constantes apropiadas $a_n$. Este es de hecho el caso, pero antes de enunciar el resultado general, consideramos un ejemplo para ilustrar la idea anterior.

\end{remark}


\begin{example}
    Determine dos soluciones en serie de potencias linealmente independientes para la ecuación diferencial
\begin{equation} \tag{9.2.2}
y'' + y = 0,
\end{equation}
centrada en $x = 0$. Identifique las soluciones en términos de funciones elementales familiares.

\noindent \textbf{Solución:} \quad Dado que $x = 0$ es un punto ordinario de la ecuación diferencial, intentamos una solución en serie de potencias de la forma
\begin{equation} \tag{4.2.3}
y(x) = \sum_{n=0}^{\infty} a_n x^n.
\end{equation}

Procedemos de manera similar al método de coeficientes indeterminados sustituyendo (4.2.3) en (4.2.2) y determinando los valores de los $a_n$ tales que (4.2.3) sea efectivamente una solución. Derivar (4.2.3) dos veces con respecto a $x$ produce
\[ y'(x) = \sum_{n=1}^{\infty} n a_n x^{n-1}, \qquad y''(x) = \sum_{n=2}^{\infty} n(n-1)a_n x^{n-2}, \]
donde hemos desplazado el punto de inicio en las sumatorias sin pérdida de generalidad. Sustituyendo en (4.2.2), se deduce que (4.2.3) define una solución, siempre que
\[ \sum_{n=2}^{\infty} n(n-1)a_n x^{n-2} + \sum_{n=0}^{\infty} a_n x^n = 0. \]

Si reemplazamos $n$ por $k + 2$ en la primera sumatoria, y reemplazamos $n$ por $k$ en la segunda sumatoria, el resultado es
\[ \sum_{k=0}^{\infty} (k+2)(k+1)a_{k+2}x^k + \sum_{k=0}^{\infty} a_k x^k = 0. \]

Combinar las sumatorias produce
\[ \sum_{k=0}^{\infty} [(k+2)(k+1)a_{k+2} + a_k]x^k = 0. \]

Esto implica que los coeficientes de $x^k$ deben anularse para $k = 0, 1, 2, \dots$. En consecuencia, obtenemos la relación de recurrencia
\[ (k+2)(k+1)a_{k+2} + a_k = 0, \qquad k = 0, 1, 2, \dots \]

Dado que $(k+2)(k+1)$ nunca es cero, podemos escribir esta relación de recurrencia en la forma equivalente
\begin{equation} \tag{4.2.4}
a_{k+2} = -\frac{1}{(k+2)(k+1)}a_k, \qquad k = 0, 1, 2, \dots
\end{equation}

Ahora usamos esta relación para determinar los valores apropiados de los coeficientes. Es conveniente considerar por separado los dos casos (1) $k$ es par y (2) $k$ es impar.

\noindent \textbf{Caso 1: $k$ es par.} \quad Sustituyendo sucesivamente en (4.2.4), obtenemos los coeficientes de la siguiente manera:
\\
Cuando $k = 0$, tenemos
$$a_2 = -\frac{1}{2}a_0.$$
Cuando $k = 2$, tenemos
$$a_4 = -\frac{1}{4 \cdot 3}a_2 = \frac{1}{4 \cdot 3 \cdot 2}a_0.$$
Es decir,
$$a_4 = \frac{1}{4!}a_0.$$
Cuando $k = 4$, tenemos
$$a_6 = -\frac{1}{6 \cdot 5}a_4,$$
de modo que
$$a_6 = -\frac{1}{6!}a_0.$$
Continuando de esta manera, pronto reconocemos el patrón que está surgiendo, a saber,
\begin{equation} \tag{4.2.5}
\boxed{ a_{2n} = \frac{(-1)^n}{(2n)!}a_0. }
\end{equation}

Así, todos los coeficientes pares están determinados en términos de $a_0$, pero $a_0$ en sí es arbitrario.

\vspace{0.5cm}
\noindent \textbf{Caso 2: $k$ es impar.} \quad Ahora considere la relación de recurrencia (4.2.4) cuando $k$ es un entero positivo impar:
\\
Cuando $k = 1$, tenemos
$$a_3 = -\frac{1}{2 \cdot 3}a_1 = -\frac{1}{3!}a_1.$$
Cuando $k = 3$, tenemos
$$a_5 = -\frac{1}{4 \cdot 5}a_3 = \frac{1}{5!}a_1.$$
Cuando $k = 5$, tenemos
$$a_7 = -\frac{1}{6 \cdot 7}a_5 = -\frac{1}{7!}a_1.$$
Continuando de esta manera, vemos que el coeficiente impar general es
\begin{equation} \tag{4.2.6}
\boxed{ a_{2n+1} = \frac{(-1)^n}{(2n+1)!}a_1. }
\end{equation}

Así, usando las Ecuaciones (4.2.5) y (4.2.6), hemos demostrado que para todos los valores de $a_0, a_1$, una solución a la ecuación diferencial dada es
$$y(x) = a_0\left(1 - \frac{1}{2!}x^2 + \frac{1}{4!}x^4 - \dots\right) + a_1\left(x - \frac{1}{3!}x^3 + \frac{1}{5!}x^5 - \dots\right).$$

Es decir,
\begin{equation} \tag{4.2.7}
y(x) = a_0 \sum_{n=0}^{\infty} \frac{(-1)^n}{(2n)!}x^{2n} + a_1 \sum_{n=0}^{\infty} \frac{(-1)^n}{(2n+1)!}x^{2n+1}.
\end{equation}

Establecer $a_0 = 1$ y $a_1 = 0$ produce la solución
$$y_1(x) = \sum_{n=0}^{\infty} \frac{(-1)^n}{(2n)!}x^{2n},$$
mientras que establecer $a_0 = 0$ y $a_1 = 1$ produce la solución
$$y_2(x) = \sum_{n=0}^{\infty} \frac{(-1)^n}{(2n+1)!}x^{2n+1}.$$

Aplicando el criterio de la razón, es sencillo demostrar que ambas series anteriores convergen para todo $x$ real. Finalmente, dado que $y_1$ y $y_2$ no son proporcionales, son linealmente independientes en cualquier intervalo. Se deduce que (4.2.7) es la solución general de la ecuación diferencial dada. De hecho, la serie de potencias que representa a $y_1$ es simplemente la expansión en serie de Maclaurin de $\cos x$, mientras que la serie que define a $y_2$ es la expansión en serie de Maclaurin de $\sin x$. Por lo tanto, podemos escribir (4.2.7) en la forma más familiar
$$y(x) = a_0 \cos x + a_1 \sin x. $$

\end{example}

El desarrollo y la soluci\'on del problema anterior, se bas\'o esencialmente en cuatro pasos, a saber: 

\begin{remark}
    \begin{enumerate}
    \item Asumir que existe una solución en serie de potencias de la forma $y(x) = \sum_{n=0}^{\infty} a_n x^n$.
    \item Determinar los valores de los coeficientes, $a_n$, de manera que $y$ sea una solución formal de la ecuación diferencial. Esto condujo a dos soluciones distintas, una determinada en términos de la constante $a_0$, y la otra en términos de la constante $a_1$.
    \item Usar el criterio de la razón para determinar el radio de convergencia de las soluciones y, por lo tanto, el intervalo sobre el cual las soluciones son válidas.
    \item Comprobar que las soluciones son linealmente independientes en el intervalo de existencia.
\end{enumerate}

\end{remark}

El siguiente teorema justifica los pasos anteriores y muestra que la técnica se puede aplicar alrededor de cualquier punto ordinario de una ecuación diferencial.

\begin{theorem}
    Sean $p$ y $q$ analíticas en $x = x_0$, y suponga que sus expansiones en serie de potencias son válidas para $|x - x_0| < R$. Entonces la solución general a la ecuación diferencial
\begin{equation} \tag{4.2.8}
y'' + p(x)y' + q(x)y = 0
\end{equation}
puede representarse como una serie de potencias centrada en $x = x_0$, con un radio de convergencia de \textit{al menos} $R$. Los coeficientes en esta solución en serie pueden determinarse en términos de $a_0$ y $a_1$ sustituyendo directamente
\[ y(x) = \sum_{n=0}^{\infty} a_n(x - x_0)^n \]
en (9.2.8). La solución resultante es de la forma
\[ y(x) = a_0 y_1(x) + a_1 y_2(x), \]
donde $y_1$ y $y_2$ son soluciones linealmente independientes de (4.2.8) en el intervalo de existencia. Si se imponen las condiciones iniciales $y(x_0) = \alpha$, $y'(x_0) = \beta$, entonces $a_0 = \alpha, a_1 = \beta$.

\end{theorem}

A continuaci\'on algunos ejemplos para darle uso al teorema anterior: 

\begin{example}
  Determine dos soluciones en serie linealmente independientes en potencias de $x$ para
\begin{equation} \tag{4.2.10}
y'' - 2xy' - 4y = 0,
\end{equation}
y encuentre el radio de convergencia de estas soluciones.

\noindent \textbf{Solución:} \quad El punto $x = 0$ es un punto ordinario de la ecuación diferencial, y por lo tanto el Teorema 4.2.4 se puede aplicar con $x_0 = 0$. En este caso, tenemos
\[ p(x) = -2x, \qquad q(x) = -4. \]

Dado que ambos son polinomios, sus expansiones en serie de potencias alrededor de $x = 0$ son válidas para todo $x$, y por lo tanto, del Teorema 4.2.4, la solución general a (4.2.10) puede expresarse en la forma
\begin{equation} \tag{4.2.11}
y(x) = \sum_{n=0}^{\infty} a_n x^n,
\end{equation}
y esta solución en serie de potencias convergerá para todo $x$ real. Derivando (4.2.11), obtenemos
\[ y'(x) = \sum_{n=1}^{\infty} n a_n x^{n-1}, \qquad y''(x) = \sum_{n=2}^{\infty} n(n-1)a_n x^{n-2}. \]
Sustituir en (4.2.10) produce
\[ \sum_{n=2}^{\infty} n(n-1)a_n x^{n-2} - 2\sum_{n=1}^{\infty} n a_n x^n - 4\sum_{n=0}^{\infty} a_n x^n = 0. \]

Ahora redefinimos los rangos en la sumatoria para obtener un $x^k$ común en todos los términos. Esto se logra reemplazando $n$ con $k + 2$ en la primera sumatoria, y, por consistencia en la notación, reemplazamos $n$ con $k$ en las otras sumatorias. El resultado es
\[ \sum_{k=0}^{\infty} [(k+2)(k+1)a_{k+2} - 2ka_k - 4a_k]x^k = 0. \]

Esta ecuación requiere que el coeficiente de $x^k$ se anule, y por lo tanto, obtenemos la relación de recurrencia
\[ (k+2)(k+1)a_{k+2} - 2ka_k - 4a_k = 0, \qquad k = 0, 1, 2, \dots, \]
que se puede escribir en la forma equivalente
\[ a_{k+2} = \frac{2(k+2)}{(k+1)(k+2)}a_k, \qquad k = 0, 1, 2, \dots; \]
es decir,
\begin{equation} \tag{4.2.12}
a_{k+2} = \frac{2}{k+1}a_k, \qquad k = 0, 1, 2, \dots
\end{equation}

Vemos por esta relación que, como en el Ejemplo 4.2.3, todos los coeficientes pares pueden expresarse en términos de $a_0$, mientras que todos los coeficientes impares pueden expresarse en términos de $a_1$. Ahora determinamos la forma exacta de estos coeficientes.

\vspace{0.3cm}
\noindent \textbf{Caso 1: $k$ es par.} \quad De (4.2.12), tenemos lo siguiente: \\
Cuando $k = 0$,
\[ a_2 = 2a_0. \]
Cuando $k = 2$,
\[ a_4 = \frac{2}{3}a_2 = \frac{2^2}{3}a_0. \]
Cuando $k = 4$,
\[ a_6 = \frac{2}{5}a_4 = \frac{2^3}{1 \cdot 3 \cdot 5}a_0. \]

El término general par es así
\[ a_{2n} = \frac{2^n}{1 \cdot 3 \cdot 5 \cdot \dots \cdot (2n - 1)}a_0. \]

\vspace{0.3cm}
\noindent \textbf{Caso 2: $k$ es impar.} \quad Sustituyendo sucesivamente en (4.2.12) se obtiene lo siguiente: \\
Cuando $k = 1$,
\[ a_3 = a_1. \]
Cuando $k = 3$,
\[ a_5 = \frac{2}{4}a_3 = \frac{1}{1 \cdot 2}a_1. \]
Cuando $k = 5$,
\[ a_7 = \frac{2}{6}a_5 = \frac{1}{1 \cdot 2 \cdot 3}a_1. \]
Cuando $k = 7$,
\[ a_9 = \frac{2}{8}a_7 = \frac{1}{1 \cdot 2 \cdot 3 \cdot 4}a_1. \]

El término general impar puede, por lo tanto, escribirse como
\[ a_{2n+1} = \frac{1}{n!}a_1. \]

Sustituyendo de nuevo en (4.2.11), obtenemos la solución
\begin{align*}
y(x) &= a_0 \left[ 1 + 2x^2 + \frac{2^2}{1 \cdot 3}x^4 + \frac{2^3}{1 \cdot 3 \cdot 5}x^6 + \dots + \frac{2^n}{1 \cdot 3 \cdot 5 \cdot \dots \cdot (2n-1)}x^{2n} + \dots \right] \\
&\quad + a_1 \left( x + x^3 + \frac{1}{2!}x^5 + \frac{1}{3!}x^7 + \dots + \frac{1}{n!}x^{2n+1} + \dots \right).
\end{align*}

Es decir,
\[ y(x) = a_0 \left[ 1 + \sum_{n=1}^{\infty} \frac{2^n}{1 \cdot 3 \cdot 5 \cdot \dots \cdot (2n-1)}x^{2n} \right] + a_1 \sum_{n=0}^{\infty} \frac{1}{n!}x^{2n+1}. \]

En consecuencia, a partir del Teorema 4.2.4, dos soluciones linealmente independientes de (4.2.10) en $(-\infty, \infty)$ son
\[ y_1(x) = 1 + \sum_{n=1}^{\infty} \frac{2^n}{1 \cdot 3 \cdot \dots \cdot (2n-1)}x^{2n}, \qquad y_2(x) = \sum_{n=0}^{\infty} \frac{1}{n!}x^{2n+1}. \tag*{$\square$} \]

\end{example}

En general, no será posible encontrar la recurrencia entre los coeficientes de las respectivas series de potencias, es por esto que en la mayor\'ia de las ocasiones, debemos ser felices con obtener una cantidad finita de términos de cada solución en serie de potencias, por ejemplo:

\begin{example}
    Determine los términos hasta $x^5$ en cada una de las dos soluciones en serie de potencias linealmente independientes para
\[ y'' + (2 - 4x^2)y' - 8xy = 0 \]
centradas en $x = 0$. Encuentre también el radio de convergencia de estas soluciones.

\noindent \textbf{Solución:} \quad Las funciones $p(x) = 2 - 4x^2$ y $q(x) = -8x$ son polinomios, y por lo tanto, a partir del Teorema 9.2.4, las soluciones en serie de potencias convergerán para todo $x$ real. Ahora determinamos las soluciones. Sustituir
\begin{equation} \tag{4.2.13}
y(x) = \sum_{n=0}^{\infty} a_n x^n
\end{equation}
en la ecuación diferencial dada produce
\[ \sum_{n=2}^{\infty} n(n-1)a_n x^{n-2} + 2\sum_{n=1}^{\infty} n a_n x^{n-1} - 4\sum_{n=1}^{\infty} n a_n x^{n+1} - 8\sum_{n=0}^{\infty} a_n x^{n+1} = 0. \]

Reemplazando $n$ por $k + 2$ en la primera sumatoria, $k + 1$ en la segunda sumatoria, y $k - 1$ en la tercera y cuarta sumatorias, obtenemos
\[ \sum_{k=0}^{\infty} (k+2)(k+1)a_{k+2}x^k + 2\sum_{k=0}^{\infty} (k+1)a_{k+1}x^k - 4\sum_{k=2}^{\infty} (k-1)a_{k-1}x^k - 8\sum_{k=1}^{\infty} a_{k-1}x^k, \]
lo cual es igual a cero. Separando los términos correspondientes a $k = 0$ y $k = 1$, se deduce que esto puede escribirse como
\begin{align*}
(2a_2 + 2a_1) + (6a_3 + 4a_2 - 8a_0)x + \quad & \\
\sum_{k=2}^{\infty} \{(k+2)(k+1)a_{k+2} + 2(k+1)a_{k+1} - [4(k-1) + 8]a_{k-1}\}x^k &= 0.
\end{align*}

Igualar los coeficientes de todas las potencias $x^k$ a cero produce lo siguiente: \\
Para $k = 0$ y $k = 1$, tenemos
\begin{equation} \tag{4.2.14}
2a_2 + 2a_1 = 0 \quad \text{y} \quad 6a_3 + 4a_2 - 8a_0 = 0,
\end{equation}
y para $k \ge 2$, tenemos
\begin{equation} \tag{4.2.15}
(k+2)(k+1)a_{k+2} + 2(k+1)a_{k+1} - [4(k-1) + 8]a_{k-1} = 0.
\end{equation}

Se deduce de (4.2.14) que
\begin{equation} \tag{4.2.16}
a_2 = -a_1, \qquad a_3 = \frac{2}{3}(2a_0 + a_1),
\end{equation}
y (4.2.15) produce la relación de recurrencia general
\begin{equation} \tag{4.2.17}
a_{k+2} = \frac{4a_{k-1} - 2a_{k+1}}{k+2}, \qquad k = 2, 3, 4, \dots
\end{equation}

En este caso, la relación de recurrencia es bastante difícil de resolver. Sin embargo, solo se nos pidió determinar los términos hasta $x^5$ en la solución en serie, por lo que procedemos a hacerlo. Ya tenemos $a_2$ y $a_3$ expresados en términos de $a_0$ y $a_1$. Establecer $k = 2$ en (4.2.17) produce
\[ a_4 = \frac{1}{4}(4a_1 - 2a_3) = \frac{1}{4}\left[ 4a_1 - \frac{4}{3}(2a_0 + a_1) \right], \]
donde hemos sustituido de (4.2.16) para $a_3$. Simplificando esta expresión, obtenemos
\[ a_4 = \frac{2}{3}(a_1 - a_0). \]

Todavía requerimos un término más. Establecer $k = 3$ en (4.2.17) produce
\[ a_5 = \frac{1}{5}(4a_2 - 2a_4) = \frac{1}{5}\left[ -4a_1 - \frac{4}{3}(a_1 - a_0) \right], \]
de modo que
\[ a_5 = \frac{4}{15}(a_0 - 4a_1). \]
Sustituyendo los coeficientes $a_2, a_3, a_4$ y $a_5$ en (4.2.13), obtenemos
\begin{align*}
y(x) &= a_0 + a_1 x - a_1 x^2 + \frac{2}{3}(2a_0 + a_1)x^3 \\
&\quad + \frac{2}{3}(a_1 - a_0)x^4 + \frac{4}{15}(a_0 - 4a_1)x^5 + \dots \ .
\end{align*}
Es decir,
\begin{align*}
y(x) &= a_0 \left( 1 + \frac{4}{3}x^3 - \frac{2}{3}x^4 + \frac{4}{15}x^5 + \dots \right) \\
&\quad + a_1 \left( x - x^2 + \frac{2}{3}x^3 + \frac{2}{3}x^4 - \frac{16}{15}x^5 + \dots \right).
\end{align*}

Así, dos soluciones linealmente independientes a la ecuación diferencial dada en $(-\infty, \infty)$ son
\[ y_1(x) = 1 + \frac{4}{3}x^3 - \frac{2}{3}x^4 + \frac{4}{15}x^5 + \dots \]
y
\[ y_2(x) = x - x^2 + \frac{2}{3}x^3 + \frac{2}{3}x^4 - \frac{16}{15}x^5 + \dots \ . \tag*{$\square$} \]

\end{example}